\documentclass[11pt,a4paper]{article}

\usepackage[utf8]{inputenc}
\usepackage[T1]{fontenc}
\usepackage{mathptmx}
\usepackage[margin=1in]{geometry}
\usepackage{amsmath,amssymb}
\usepackage{booktabs}
\usepackage{array}
\usepackage{hyperref}
\usepackage{float}
\usepackage{xcolor}
\usepackage{natbib}

\hypersetup{
    colorlinks=true,
    linkcolor=blue,
    citecolor=blue,
    urlcolor=blue,
}

\title{\LARGE\textbf{Supplementary Material}\\
\large Mathematical Definitions, Notation, and Computational Procedures\\
\normalsize Paper 6: A Conservation Constraint in LLM Context Processing\\
Across Four Epistemological Domains}

\author{Dr.\ Laxman M M, MBBS\\
Government Duty Medical Officer, PHC Manchi\\
Bantwal Taluk, Dakshina Kannada, Karnataka, India}

\date{April 2026}

\begin{document}

\maketitle
\tableofcontents
\newpage

% ================================================================
\section{Experimental Conditions}

Each model-domain run consists of three conditions, each comprising 50 independent
trials of a 30-prompt conversation.

\begin{description}
    \item[TRUE condition] Each trial presents the model with a coherent 29-message
    conversational history followed by a prompt. The history is the same across all
    50 trials within a run. Context is semantically and temporally coherent.

    \item[COLD condition] Each trial presents the model with the same prompt in
    isolation, with no prior conversational history. This is the context-free baseline.

    \item[SCRAMBLED condition] Each trial presents the model with the same 29 messages
    from the TRUE history, but in a uniformly randomised order, followed by the same
    prompt. Content is present; temporal order is destroyed.
\end{description}

\noindent All conditions use temperature $= 0.7$ and max\_tokens $= 1024$.
Embeddings are computed using all-MiniLM-L6-v2 (384-dimensional) as the primary model,
with all-mpnet-base-v2 (768-dimensional) and LaBSE (768-dimensional) for robustness
validation.

% ================================================================
\section{Core Metrics}

\subsection{Relational Coherence Index (RCI)}

For a given condition $c \in \{\text{TRUE}, \text{COLD}, \text{SCRAMBLED}\}$,
model $m$, domain $d$, and prompt position $p \in \{1, \ldots, 30\}$:

Let $\mathbf{e}_{i,p}^{(c)}$ denote the embedding of the response generated in trial
$i$ at position $p$ under condition $c$.

The RCI at position $p$ is the mean pairwise cosine similarity across $N = 50$ trials:

\begin{equation}
    \text{RCI}_{p}^{(c)} = \frac{2}{N(N-1)} \sum_{i < j}
    \cos\!\big(\mathbf{e}_{i,p}^{(c)},\, \mathbf{e}_{j,p}^{(c)}\big)
\end{equation}

The model-level RCI is the mean across all 30 positions:

\begin{equation}
    \overline{\text{RCI}}^{(c)} = \frac{1}{30} \sum_{p=1}^{30} \text{RCI}_{p}^{(c)}
\end{equation}

\subsection{Delta Relational Coherence Index ($\Delta$RCI)}

$\Delta$RCI measures how much the conversational context increases response coherence
relative to the context-free baseline:

\begin{equation}
    \Delta\text{RCI} = \overline{\text{RCI}}^{(\text{TRUE})} -
    \overline{\text{RCI}}^{(\text{COLD})}
\end{equation}

$\Delta\text{RCI} > 0$ indicates that context increases response alignment across
trials. $\Delta\text{RCI} < 0$ indicates divergent context coupling (context
\emph{increases} response variability).

\subsection{Variance Ratio (Var\_Ratio)}

The variance of responses at position $p$ under condition $c$ is computed from the
pairwise cosine similarity distribution across all $N = 50$ trial pairs. Let:

\begin{equation}
    S_p^{(c)} = \left\{ \cos\!\big(\mathbf{e}_{i,p}^{(c)},\, \mathbf{e}_{j,p}^{(c)}\big)
    : 1 \leq i < j \leq N \right\}
\end{equation}

be the set of all $\binom{N}{2} = 1225$ pairwise cosine similarities at position $p$.
The positional variance is:

\begin{equation}
    \text{Var}_{p}^{(c)} = \text{Var}\!\left( S_p^{(c)} \right)
\end{equation}

The model-level Var\_Ratio is:

\begin{equation}
    \text{Var\_Ratio} = \frac{\overline{\text{Var}}^{(\text{TRUE})}}
    {\overline{\text{Var}}^{(\text{COLD})}}
    \quad \text{where} \quad
    \overline{\text{Var}}^{(c)} = \frac{1}{30} \sum_{p=1}^{30} \text{Var}_{p}^{(c)}
\end{equation}

This formulation measures the spread of the pairwise similarity distribution rather
than dimension-wise embedding variance, capturing how consistently the model produces
similar responses across trials at each position.

$\text{Var\_Ratio} < 1$: context reduces response variance (convergent coupling).\\
$\text{Var\_Ratio} > 1$: context increases response variance (divergent coupling).

\subsection{Variance Reduction Index (VRI)}

\begin{equation}
    \text{VRI} = 1 - \text{Var\_Ratio}
\end{equation}

VRI is used in entanglement analysis (Section~\ref{sec:entanglement}). Positive VRI
indicates convergent coupling; negative VRI indicates divergent coupling.

% ================================================================
\section{The Conservation Product K}

\subsection{Definition}

For a given model-domain run, the conservation product is:

\begin{equation}
    K = \Delta\text{RCI} \times \text{Var\_Ratio}
\end{equation}

\subsection{Domain-Level Statistics}

For domain $d$ with $N_d$ model-domain runs yielding values
$K_1^{(d)}, K_2^{(d)}, \ldots, K_{N_d}^{(d)}$:

\begin{align}
    \bar{K}^{(d)} &= \frac{1}{N_d} \sum_{i=1}^{N_d} K_i^{(d)}
    \quad \text{(domain mean)} \\[6pt]
    \text{SD}^{(d)} &= \sqrt{\frac{1}{N_d - 1} \sum_{i=1}^{N_d}
    \left(K_i^{(d)} - \bar{K}^{(d)}\right)^2}
    \quad \text{(standard deviation)} \\[6pt]
    \text{CV}^{(d)} &= \frac{\text{SD}^{(d)}}{\bar{K}^{(d)}}
    \quad \text{(coefficient of variation)}
\end{align}

CV is used as the primary measure of within-domain conservation tightness.
Lower CV indicates tighter clustering of $K$ values within the domain.

\subsection{Domain Values (Primary Embedding: MiniLM 384D)}

\begin{table}[H]
\centering
\begin{tabular}{lcccc}
\toprule
\textbf{Domain} & \textbf{N} & $\bar{K}$ & \textbf{SD} & \textbf{CV} \\
\midrule
Medical (Discovered)     & 8 & 0.429 & 0.073 & 0.170 \\
Legal (Argued)           & 5 & 0.348 & 0.074 & 0.214 \\
Philosophy (Explored)    & 6 & 0.301 & 0.050 & 0.166 \\
Applied Ethics (Felt)    & 5 & 0.223 & 0.036 & 0.162 \\
\bottomrule
\end{tabular}
\caption{Domain-level $K$ statistics under primary embedding (all-MiniLM-L6-v2).}
\end{table}

% ================================================================
\section{Content-Order Decomposition}
\label{sec:content_order}

The SCRAMBLED condition separates $\Delta$RCI into two components:

\begin{align}
    \Delta\text{RCI}_{\text{content}} &= \overline{\text{RCI}}^{(\text{SCRAM})} -
    \overline{\text{RCI}}^{(\text{COLD})}
    \quad \text{(content component)} \\[4pt]
    \Delta\text{RCI}_{\text{order}} &= \overline{\text{RCI}}^{(\text{TRUE})} -
    \overline{\text{RCI}}^{(\text{SCRAM})}
    \quad \text{(order component)}
\end{align}

The content fraction is:

\begin{equation}
    \text{Content Fraction} = \frac{\Delta\text{RCI}_{\text{content}}}
    {\Delta\text{RCI}} \times 100\%
    = \frac{\overline{\text{RCI}}^{(\text{SCRAM})} - \overline{\text{RCI}}^{(\text{COLD})}}
    {\overline{\text{RCI}}^{(\text{TRUE})} - \overline{\text{RCI}}^{(\text{COLD})}}
    \times 100\%
\end{equation}

Content Fraction $= 100\%$ means SCRAMBLED $\approx$ TRUE: conversational order
carries no additional information. Content Fraction $= 0\%$ means SCRAMBLED $\approx$
COLD: all context sensitivity is carried by order, not content.

% ================================================================
\section{Exploration Arc}
\label{sec:arc}

The Exploration Arc measures whether response variance converges or diverges over the
course of the 30-prompt conversation:

\begin{equation}
    \text{Arc} = \frac{\text{Var}_{30}^{(\text{TRUE})}}{\text{Var}_{1}^{(\text{TRUE})}}
\end{equation}

Arc $< 1$: responses converge (become more similar across trials) as conversation
deepens. This is the pattern for Medical (Discovered truth), where diagnostic
conclusions narrow with accumulated context.

Arc $> 1$: responses diverge (become more variable across trials) as conversation
deepens. This is the pattern for Philosophy (Explored truth), where inquiry expands
rather than converges.

% ================================================================
\section{Entanglement}
\label{sec:entanglement}

Entanglement measures the coupling between context sensitivity and variance reduction
across the 30 prompt positions within a single model-domain run.

For each position $p$, compute position-level $\Delta\text{RCI}_p$ and
$\text{VRI}_p = 1 - \text{Var\_Ratio}_p$.

Entanglement is the Pearson correlation across all 30 positions:

\begin{equation}
    r_{\text{ent}} = \text{Pearson}\!\left(
    \{\Delta\text{RCI}_p\}_{p=1}^{30},\;
    \{\text{VRI}_p\}_{p=1}^{30}
    \right)
\end{equation}

Statistical significance tested at $\alpha = 0.05$ (two-sided).

Positive $r_{\text{ent}}$: positions where context sensitivity is high also show
greater variance reduction. This is \emph{convergent entanglement}.

Negative $r_{\text{ent}}$: positions where context sensitivity is high show greater
variance \emph{amplification}. This is \emph{divergent entanglement}.

% ================================================================
\section{Effective Rank (from Paper 9)}

The effective rank of an embedding matrix $\mathbf{E} \in \mathbb{R}^{N \times D}$
is computed as the exponential of the Shannon entropy of the normalised singular
value spectrum:

\begin{equation}
    \text{eff\_rank}(\mathbf{E}) = \exp\!\left(H(\mathbf{p})\right)
    \quad \text{where} \quad
    p_k = \frac{\sigma_k}{\sum_j \sigma_j}
\end{equation}

and $\sigma_1 \geq \sigma_2 \geq \cdots \geq \sigma_{\min(N,D)} \geq 0$ are the
singular values of $\mathbf{E}$ and $H(\mathbf{p}) = -\sum_k p_k \log p_k$ is the
Shannon entropy of the normalised spectrum.

Effective rank $\approx \min(N,D)$ indicates a full-rank, non-degenerate embedding
space. Effective rank $\approx 1$ indicates near-complete collapse to a single
dimension. MuRIL achieves effective rank $= 1.42/10$; MiniLM achieves $9.88/10$
(Paper 9, \url{https://doi.org/10.5281/zenodo.19466613}).

% ================================================================
\section{Entanglement Stability Index (ESI, from Paper 4)}
\label{sec:esi}

The Entanglement Stability Index predicts which models will exhibit
multi-turn instability (the ``Lost in Conversation'' phenomenon):

\begin{equation}
    \text{ESI} = \frac{1}{\left| 1 - \text{Var\_Ratio} \right|}
\end{equation}

ESI measures the proximity of Var\_Ratio to unity. When
$\text{Var\_Ratio} \approx 1$, context neither amplifies nor
reduces variance, yielding high ESI (stable). When Var\_Ratio
deviates substantially from 1 (either convergent or divergent),
ESI is low (unstable).

\begin{itemize}
    \item ESI $> 2.0$: Stable multi-turn behaviour
    \item ESI $< 1.0$: Elevated instability risk
    \item ESI $\to \infty$: Var\_Ratio $= 1.0$ (perfect stability)
    \item ESI $\to 0$: Extreme divergent or convergent coupling
\end{itemize}

Paper 4 \citep{laxman2026d} demonstrated that ESI $< 1.0$ at task-critical
positions (Medical P30) identifies models exhibiting the ``Lost in Conversation''
phenomenon documented by Laban et al.\ (2025) across 200,000+ conversations.

\begin{table}[H]
\centering
\begin{tabular}{lccc}
\toprule
\textbf{Model} & \textbf{Var\_Ratio (P30)} & \textbf{ESI} & \textbf{Assessment} \\
\midrule
Llama 4 Scout    & 7.46 & 0.15 & Unstable \\
Llama 4 Maverick & 2.64 & 0.61 & Caution  \\
Qwen3 235B       & 1.02 & 50.0 & Stable   \\
GPT-4o           & 0.58 & 2.38 & Stable   \\
Claude Haiku     & 0.65 & 2.86 & Stable   \\
\bottomrule
\end{tabular}
\caption{ESI at medical summarisation position (P30). From Paper 4
\citep{laxman2026d}.}
\end{table}

% ================================================================
\section{Encoding Fidelity Index (EFI)}
\label{sec:efi}

Introduced in Paper 8. EFI measures how faithfully an LLM encodes non-English input
relative to its English equivalent.

\subsection{Theoretical Definition}

EFI is defined as the normalised mutual information between the source message $X$
and the model's internal representation $\hat{X}$:

\begin{equation}
    \text{EFI} = \frac{I(X;\, \hat{X})}{H(X)}
\end{equation}

where $H(X)$ is the Shannon entropy of the source. EFI $= 1$ indicates perfect
encoding fidelity; EFI $= 0$ indicates complete encoding failure.

Since $I(X; \hat{X})$ requires access to model internals, a practical proxy is used.

\subsection{Proxy Definition}

\begin{equation}
    \text{EFI}_{\text{proxy}} = \cos\!\big(\mathbf{e}(\text{non-English input}),\,
    \mathbf{e}(\text{English equivalent})\big)
\end{equation}

where $\mathbf{e}(\cdot)$ is a sentence embedding. English self-similarity serves as
the reference: $\text{EFI}_{\text{proxy}}(\text{English}) = 1.0$ by construction.

\subsection{Measurement Instrument Dependence}

Paper 9 established that EFI$_{\text{proxy}}$ is substantially instrument-dependent.
Under the primary embedding (MiniLM), Kannada EFI $= 0.081$ --- a 92\% degradation.
Under LaBSE (language-agnostic), Kannada EFI $= 0.853$ --- only 15\% below the
English reference. The phenomenon is real under all instruments; its measured
magnitude is conservative under English-distilled models.

\begin{table}[H]
\centering
\begin{tabular}{lcc}
\toprule
\textbf{Language} & \textbf{EFI (MiniLM)} & \textbf{EFI (LaBSE)} \\
\midrule
English  & 1.000 & 1.000 \\
Kannada  & 0.081 & 0.853 \\
Tamil    & 0.073 & 0.857 \\
Hindi    & 0.041 & 0.861 \\
\bottomrule
\end{tabular}
\caption{EFI by language and embedding model (from Papers 8 and 9).}
\end{table}

\subsection{Relationship to Var\_Ratio}

EFI and Var\_Ratio are statistically independent under LaBSE ($r = -0.18$,
$p = 0.73$, $N = 6$; Paper 9). They measure distinct phenomena: EFI measures
encoding quality at the input stage; Var\_Ratio measures output consistency at the
generation stage. Both are components of the MCH safety framework but require
separate evaluation.

% ================================================================
\section{Statistical Tests}

\subsection{Domain Separation}

Pairwise domain separation is tested using the Mann-Whitney $U$ statistic
(non-parametric, two-sided). Effect sizes reported as Cohen's $d$:

\begin{equation}
    d = \frac{\bar{K}^{(d_1)} - \bar{K}^{(d_2)}}
    {\sqrt{\left(\text{SD}^{(d_1)2} + \text{SD}^{(d_2)2}\right)/2}}
\end{equation}

\subsection{Entanglement Significance}

Pearson $r$ tested against $H_0: r = 0$ using the $t$-distribution with $df = 28$
(two-sided). Significance codes: * $p < 0.05$, ** $p < 0.01$, *** $p < 0.001$.

\subsection{Bootstrap Confidence Intervals (from Paper 9)}

For Var\_Ratio confidence intervals: 10,000 bootstrap iterations, resampling
trials with replacement. 95\% CI lower bound exceeding 1.0 indicates significant
variance amplification.

% ================================================================
\section{Embedding Model Specifications}

\begin{table}[H]
\centering
\begin{tabular}{llll}
\toprule
\textbf{Model} & \textbf{Dim} & \textbf{Category} & \textbf{Role in Paper 6} \\
\midrule
all-MiniLM-L6-v2  & 384 & EN-distilled  & Primary embedding \\
all-mpnet-base-v2 & 768 & EN-distilled  & Robustness check 1 \\
LaBSE             & 768 & Language-agnostic & Robustness check 2 \\
\bottomrule
\end{tabular}
\caption{Embedding models used for $K$ computation. LaBSE validated as semantically
faithful in Paper 9 (non-parallel control: $\Delta = 0.58$, $p = 0.004$).}
\end{table}

% ================================================================
\section{Relationship Between Metrics}

The following relationships hold by construction or empirical observation:

\begin{align}
    K &= \Delta\text{RCI} \times \text{Var\_Ratio} \\
    \text{VRI} &= 1 - \text{Var\_Ratio} \\
    \Delta\text{RCI} &= \Delta\text{RCI}_{\text{content}} +
    \Delta\text{RCI}_{\text{order}} \\
    r_{\text{ent}} &\approx +0.76 \text{ in Medical domain (Paper 4)}
\end{align}

The conservation constraint states that $K$ is approximately constant within a
domain across model architectures:

\begin{equation}
    K_i^{(d)} \approx \bar{K}^{(d)} \quad \forall\, i \in \{1, \ldots, N_d\}
\end{equation}

with CV $= 0.162$--$0.214$ under MiniLM, indicating approximate rather than exact
conservation.

% ================================================================
\section{Data and Code Availability}

All raw data, analysis scripts, and computed metrics are available at:
\begin{itemize}
    \item GitHub: \url{https://github.com/LaxmanNandi/MCH-Research}
    \item OSF pre-registration: \url{https://osf.io/dp8nj/}
\end{itemize}

Key scripts:
\begin{itemize}
    \item \texttt{paper6\_conservation.py} --- K computation across all domains
    \item \texttt{paper6\_robustness\_reembed.py} --- Three-embedding robustness
    \item \texttt{paper6\_entanglement.py} --- Entanglement analysis
    \item \texttt{compile\_paper6\_data.py} --- Final data compilation to JSON
\end{itemize}

\bibliographystyle{plainnat}
\begin{thebibliography}{5}

\bibitem[Laxman(2026d)]{laxman2026d}
Laxman, M.~M. (2026d).
\newblock Engagement as entanglement: Variance signatures of bidirectional
context coupling in large language models.
\newblock \emph{Preprints.org}.
\newblock \url{https://doi.org/10.20944/preprints202603.0055.v1}

\bibitem[Laxman(2026h)]{laxman2026h}
Laxman, M.~M. (2026h).
\newblock Encoding fidelity and coherent misalignment in non-{E}nglish clinical {AI}.
\newblock \emph{Preprints.org}.
\newblock \url{https://doi.org/10.20944/preprints202604.0061.v1}

\bibitem[Laxman(2026i)]{laxman2026i}
Laxman, M.~M. (2026i).
\newblock Measurement matters: Embedding model choice determines encoding
fidelity assessment in multilingual clinical {AI}.
\newblock \emph{Zenodo}.
\newblock \url{https://doi.org/10.5281/zenodo.19466613}

\end{thebibliography}

\end{document}
